The QSPI memory was tested this version as I discovered the rudimentary testing I completed for v0.1 had flaws and errors e.g., misinterpreting the prescalars. These tests aimed to provide an understanding of how long it takes to complete transfers, which helped in determining system speeds for the camera.
\nl
To observe how long it took to complete a write transfer, I first verified that the write command was completing successfully via several 1 MB write/read cycles which revealed no errors. I then tracked how many system ticks between starting the write command and finishing the command, with reading and verification occuring afterwards. With different prescalars and kernel speeds, I plotted the data to observe the relationship between peripheral speed and data transfer in \autoref{fig:v1-qspi-data}.
\nl
The time constant represents approximately how long it takes to transmit $x$ bytes in nanoseconds. I found that, as expected, the time constant reduces in the $1/x$ behaviour we would expect from a system like this which proved: (a) faster speeds transmit faster and; (b) there is a predictable relationship between the frequency and time constant. Moreover, I was able to verify the reciprocal relationship with measurements from a logic analyser in \autoref{fig:v1-qspi-prescalar}.
\nl
The most notable takeaway from the testing was that as the clock frequency increased, the time constant assymptoted towards a value of $2220$ ns/B. This indicated the physical barrier of the system i.e., how fast it can possibly go. This value is what was used in \autoref{eq:ov7670-pclk} to determine the camera prescalar value.
\nl
I found that optimising the QSPI driver, by removing safety checks, reduced this time constant (higher data rate), which suggested that this limit wasn't entirely generated by hardware restrictions, but code as well. As \href{https://www.reddit.com/r/embedded/comments/1v77w8s/comment/p04h2kg/?screen_view_count=2&ext-referrer=DIRECT}{Neywiny\_} outlined, the timeline is proportional to $C_1 + C_2 D + D/C_3$, where $C_1$ represents the HAL (hardware limitation), $C_3$ the clock frequency, $C_2$ is the system link, and $D$ the data size. We can see that this experiment varied $C_3$ and proved that faster clock speeds and smaller data payloads reduce the transfer time. Although the commenter mentioned $C_1$ could be removed by using DMA, we were already using the GPDMA for the DCMI request, so that was not a possible approach.
\nl
Overall, this basic testing helped characterise the time constant of the current QSPI implementation to aid with system configuration.

\begin{figure}
    \centering
    \includegraphics[width=0.7\textwidth]{Images/v1/Testing/QSPI.png}
    \caption{Plot of time constant against effective QSPI clock speed.}
    \label{fig:v1-qspi-data}
\end{figure}

\begin{figure}
    \centering
    \includegraphics[width=0.7\textwidth]{Images/v1/Testing/QSPI-Prescalar.png}
    \caption{Comparison between QSPI CLK speed of predicted and measured speeds.}
    \label{fig:v1-qspi-prescalar}
\end{figure}